\section{The Introduction}
\label{sec:write}

Advice on writing introductions is usually folklore; this section
states it as a grammar.
We reverse engineered four introductions, nine decades apart, for
the shape they share: Shannon's thesis on switching
circuits~\cite{shannon40}, an effort-estimation study~\cite{kocaguneli12},
a deep-learning replication~\cite{fu17}, and a security
taxonomy~\cite{ladisa23}.
Widom's writing guide supplied the audit questions~\cite{widom06}.
Figure~\ref{fig:introbnf} shows the result as a definite clause
grammar.

The notation, which the figure borrows from logic programming,
reads as follows.
\begin{itemize}
\item Lower-case names are nonterminals; {\ttfamily -->} rewrites
the left side into the sequence on the right, and commas order that
sequence.
\item Semicolons separate alternatives; hence a hook is any one of
its three forms.
\item Bracketed names are leaves, and each leaf is one prose move,
a sentence to a short paragraph long.
\item {\ttfamily +} means one or more; {\ttfamily *} means zero or
more.
\end{itemize}

Each production also carries a rule of craft, as follows.
\begin{itemize}
\item {\ttfamily hook} opens with a number that hurts (dollars,
hours, incident counts), an anchor to a prior result, or a lived
setting.
The importance argument lives inside those numbers, which is why no
separate why-this-matters clause exists.
\item {\ttfamily gap} walks the Stanford four clauses in order.
Name and define the problem.
Show why the naive approach fails, and why prior work has not
already fixed it.
State the key idea in one sentence, which the rest of the paper
unpacks.
\item {\ttfamily elevator} states the hypothesis as the conclusion,
displayed as a quote of two to three lines.
Papers carry at most one or two displayed claims, and the elevator
speech spends one of them.
\item {\ttfamily move} says what this paper does, then christens
it; the method gets a name at first use and keeps it ever after.
\item {\ttfamily rqs} answers every question on the spot, its
headline number in bold, since an introduction holds no suspense.
\item {\ttfamily contribs} is a short bullet list whose final
bullet is always the replication URL.
\item {\ttfamily [caveat]*} allows zero or more digressions, each
answering an obvious objection before a reviewer raises it.
\end{itemize}

\begin{pause}{the introduction}
Ask: would your own last paper parse under this grammar, and which
leaf would fail first?
Does your field replace any production (some venues demand a
roadmap paragraph; this grammar omits one)?\\[0.2em]
To change the outcome: edit the grammar here, then regenerate the
sample of Figure~\ref{fig:sample} against the new productions.
\end{pause}
